% Chunk: \S5.5 Causal Dirac Fields
% Source: Weinberg QFT Vol. I, physical pp. 242--252 (printed pp. 219--229; begins and ends mid-page)
% Status: source-reviewed, compile-clean, and visually verified; equations (5.5.1)--(5.5.58); no figures; no uncertainties

\subsection{Causal Dirac Fields}

We now want to construct particle annihilation and antiparticle creation fields
that transform under the Lorentz group according to the Dirac representation
of this group, discussed in the previous section. In general these take the
form given in Eqs.~\eqref{eq:5.1.17} and \eqref{eq:5.1.18}:
\begin{equation}
\psi_\ell^{(+)}(x)=(2\pi)^{-3/2}\sum_\sigma\int d^3p\,
u_\ell(\mathbf p,\sigma)e^{ip\cdot x}a_{\mathbf p,\sigma}
\tag{5.5.1}\label{eq:5.5.1}
\end{equation}
and
\begin{equation}
\psi_\ell^{(-c)}(x)=(2\pi)^{-3/2}\sum_\sigma\int d^3p\,
v_\ell(\mathbf p,\sigma)e^{-ip\cdot x}a^{c\dagger}_{\mathbf p,\sigma},
\tag{5.5.2}\label{eq:5.5.2}
\end{equation}
with the particle species label omitted here. In order to calculate the
coefficient functions $u_\ell(\mathbf p,\sigma)$ and
$v_\ell(\mathbf p,\sigma)$ appearing in these formulas, we must first use Eqs.
\eqref{eq:5.1.25} and \eqref{eq:5.1.26} to find $u_\ell$ and $v_\ell$ for zero momentum, and then
apply Eqs.~\eqref{eq:5.1.21} and \eqref{eq:5.1.22} to calculate them for arbitrary momenta, with
$D_{\ell\bar\ell}(\Lambda)$ in both cases taken as the $4\cdot4$ Dirac
representation of the homogeneous Lorentz group discussed in the previous
section.

Using Eq.~\eqref{eq:5.4.19}, the zero-momentum conditions \eqref{eq:5.1.25} and \eqref{eq:5.1.26}
read\footnote{We are here dropping the species label $n$, and replacing the
four-component index $\ell$ with a pair of indices, one 2-valued index $m$
labelling the rows and columns of the submatrices in Eqs.~\eqref{eq:5.4.19} and
\eqref{eq:5.4.20}, and a second index taking values $\mathord\pm$, labelling the rows
and columns of the supermatrix in Eqs.~\eqref{eq:5.4.19} and \eqref{eq:5.4.20}.}
\[
\sum_{\bar\sigma}u_{m,\pm}(\mathbf0,\bar\sigma)
J^{(j)}_{\bar\sigma\sigma}
=\sum_{\bar m}\frac12\boldsymbol{\sigma}_{m\bar m}
u_{\bar m,\pm}(\mathbf0,\sigma)
\]
and
\[
-\sum_{\bar\sigma}v_{m,\pm}(\mathbf0,\bar\sigma)
J^{(j)*}_{\bar\sigma\sigma}
=\sum_{\bar m}\frac12\boldsymbol{\sigma}_{m\bar m}
v_{\bar m,\pm}(\mathbf0,\sigma).
\]
In other words, if we regard $u_{m,\pm}(\mathbf0,\sigma)$ and
$v_{m,\pm}(\mathbf0,\sigma)$ as the $m,\sigma$ elements of matrices $U_\pm$
and $V_\pm$, we have in matrix notation
\begin{equation}
U_\pm\mathbf J^{(j)}=\frac12\boldsymbol{\sigma}U_\pm
\tag{5.5.3}\label{eq:5.5.3}
\end{equation}
and
\begin{equation}
-V_\pm\mathbf J^{(j)*}=\frac12\boldsymbol{\sigma}V_\pm.
\tag{5.5.4}\label{eq:5.5.4}
\end{equation}

Now, the $(2j+1)$-dimensional matrices $\mathbf J^{(j)}$ and
$-\mathbf J^{(j)*}$ and the $2\cdot2$ matrices
$\frac12\boldsymbol{\sigma}$ all provide irreducible representations of the
Lie algebra of the rotation group. A general theorem of group theory known as
Schur's lemma~\hyperref[ref:5.6]{[6]} tells us that when a matrix like $U_\pm$ or $V_\pm$ connects
two such representations as in Eqs.~\eqref{eq:5.5.3} and \eqref{eq:5.5.4}, the matrix must
either vanish (a possibility of no interest here) or else be square and
non-singular. Hence the Dirac field can only describe particles of spin
$j=\frac12$ (so that $2j+1=2$) and the matrices $\mathbf J^{(1/2)}$ and
$-\mathbf J^{(1/2)*}$ must be the same as
$\frac12\boldsymbol{\sigma}$ up to a similarity transformation. In fact, in
the standard representation \eqref{eq:2.5.21}, \eqref{eq:2.5.22} of the rotation generators, we
have $\mathbf J^{(1/2)}=\frac12\boldsymbol{\sigma}$ and
$-\mathbf J^{(1/2)*}=\frac12\sigma_2\boldsymbol{\sigma}\sigma_2$. It follows
then that $U_\pm$ and $V_\pm\sigma_2$ must commute with
$\boldsymbol{\sigma}$, and hence must be proportional to the unit matrix:
\begin{equation}
u_{m,\pm}(\mathbf0,\sigma)=c_\pm\delta_{m\sigma},
\qquad
v_{m,\pm}(\mathbf0,\sigma)=i d_\pm(\sigma_2)_{m\sigma}.
\tag{5.5.5}\label{eq:5.5.5}
\end{equation}
In other words
\[
u(\mathbf0,\tfrac12)=\begin{bmatrix}c_+\\0\\c_-\\0\end{bmatrix},
\quad
u(\mathbf0,-\tfrac12)=\begin{bmatrix}0\\c_+\\0\\c_-\end{bmatrix},
\]
\[
v(\mathbf0,\tfrac12)=\begin{bmatrix}0\\d_+\\0\\d_-\end{bmatrix},
\quad
v(\mathbf0,-\tfrac12)=-\begin{bmatrix}d_+\\0\\d_-\\0\end{bmatrix},
\]
and the spinors at finite momentum are
\begin{equation}
u(\mathbf p,\sigma)=\sqrt{m/p^0}\,D(L(\mathbf p))u(\mathbf0,\sigma),
\tag{5.5.6}\label{eq:5.5.6}
\end{equation}
\begin{equation}
v(\mathbf p,\sigma)=\sqrt{m/p^0}\,D(L(\mathbf p))v(\mathbf0,\sigma).
\tag{5.5.7}\label{eq:5.5.7}
\end{equation}
It now only remains to say something about the constants $c_\pm$ and $d_\pm$.
In general, these are quite arbitrary---we could even choose $c_-$ and $d_-$
or $c_+$ and $d_+$ to be zero if we liked, so that the Dirac field would have
only two non-vanishing components. The only physical principle that could tell
us anything about the relative values of the $c_\pm$ or the $d_\pm$ is the
conservation of parity. We recall that under a space inversion, the particle
annihilation and antiparticle creation operators undergo the transformations:
\begin{equation}
P a_{\mathbf p,\sigma}P^{-1}=\eta^*a_{-\mathbf p,\sigma}
\tag{5.5.8}\label{eq:5.5.8}
\end{equation}
\begin{equation}
P a^{c\dagger}_{\mathbf p,\sigma}P^{-1}
=\eta^c a^{c\dagger}_{-\mathbf p,\sigma}
\tag{5.5.9}\label{eq:5.5.9}
\end{equation}
and so
\begin{equation}
P\psi_\ell^{(+)}(x)P^{-1}=\eta^*(2\pi)^{-3/2}
\sum_\sigma\int d^3p\,u_\ell(-\mathbf p,\sigma)
e^{ip\cdot\mathcal P x}a_{\mathbf p,\sigma},
\tag{5.5.10}\label{eq:5.5.10}
\end{equation}
\begin{equation}
P\psi_\ell^{(-c)}(x)P^{-1}=\eta^c(2\pi)^{-3/2}
\sum_\sigma\int d^3p\,v_\ell(-\mathbf p,\sigma)
e^{-ip\cdot\mathcal P x}a^{c\dagger}_{\mathbf p,\sigma}.
\tag{5.5.11}\label{eq:5.5.11}
\end{equation}
Also, Eqs.~\eqref{eq:5.4.16}, \eqref{eq:5.1.21}, and \eqref{eq:5.1.22} give
\begin{equation}
u(-\mathbf p,\sigma)=\sqrt{m/p^0}\,\beta D(L(\mathbf p))\beta u(\mathbf0,\sigma)
\tag{5.5.12}\label{eq:5.5.12}
\end{equation}
\begin{equation}
v(-\mathbf p,\sigma)=\sqrt{m/p^0}\,\beta D(L(\mathbf p))\beta v(\mathbf0,\sigma).
\tag{5.5.13}\label{eq:5.5.13}
\end{equation}
(Since $\beta^2=\mathbf{1}$, we are no longer making a distinction between
$\beta$ and $\beta^{-1}$.) In order that the parity operator should transform
the annihilation and creation fields at the point $x$ into something
proportional to these fields at $\mathcal P x$, it is necessary that
$\beta u(\mathbf0,\sigma)$ and $\beta v(\mathbf0,\sigma)$ be proportional to
$u(\mathbf0,\sigma)$ and $v(\mathbf0,\sigma)$, respectively:
\begin{equation}
\beta u(\mathbf0,\sigma)=b_u u(\mathbf0,\sigma),
\qquad
\beta v(\mathbf0,\sigma)=b_v v(\mathbf0,\sigma),
\tag{5.5.14}\label{eq:5.5.14}
\end{equation}
where $b_u$ and $b_v$ are sign factors, $b_u^2=b_v^2=1$. In this case, the
fields have the simple space-inversion properties:
\begin{equation}
P\psi^{(+)}(x)P^{-1}=\eta^*b_u\beta\psi^{(+)}(\mathcal P x),
\tag{5.5.15}\label{eq:5.5.15}
\end{equation}
\begin{equation}
P\psi^{(-c)}(x)P^{-1}=\eta^c b_v\beta\psi^{(-c)}(\mathcal P x).
\tag{5.5.16}\label{eq:5.5.16}
\end{equation}
By adjusting the overall scales of the fields, we can choose the coefficient
functions at zero momentum to have the form:
\begin{equation}
u(\mathbf0,\tfrac12)=\frac1{\sqrt2}\begin{bmatrix}1\\0\\b_u\\0\end{bmatrix},
\qquad
u(\mathbf0,-\tfrac12)=\frac1{\sqrt2}\begin{bmatrix}0\\1\\0\\b_u\end{bmatrix},
\tag{5.5.17}\label{eq:5.5.17}
\end{equation}
\begin{equation}
v(\mathbf0,\tfrac12)=\frac1{\sqrt2}\begin{bmatrix}0\\1\\0\\b_v\end{bmatrix},
\qquad
v(\mathbf0,-\tfrac12)=-\frac1{\sqrt2}\begin{bmatrix}1\\0\\b_v\\0\end{bmatrix}.
\tag{5.5.18}\label{eq:5.5.18}
\end{equation}

Now let's try to put together the annihilation and creation fields in a linear
combination
\begin{equation}
\psi(x)=\kappa\psi^{(+)}(x)+\lambda\psi^{(-c)}(x)
\tag{5.5.19}\label{eq:5.5.19}
\end{equation}
that commutes or anticommutes with itself and its adjoint at space-like
separations. A straightforward calculation gives
\begin{align}
[\psi_\ell(x),\psi_{\bar\ell}^\dagger(y)]_\mp
=\int d^3p\bigl[|\kappa|^2N_{\ell\bar\ell}(\mathbf p)e^{ip\cdot(x-y)}
\mp|\lambda|^2M_{\ell\bar\ell}(\mathbf p)e^{-ip\cdot(x-y)}\bigr],
\tag{5.5.20}\label{eq:5.5.20}
\end{align}
where
\begin{equation}
N_{\ell\bar\ell}(\mathbf p)\equiv
\sum_\sigma u_\ell(\mathbf p,\sigma)u_{\bar\ell}^*(\mathbf p,\sigma),
\tag{5.5.21}\label{eq:5.5.21}
\end{equation}
\begin{equation}
M_{\ell\bar\ell}(\mathbf p)\equiv
\sum_\sigma v_\ell(\mathbf p,\sigma)v_{\bar\ell}^*(\mathbf p,\sigma).
\tag{5.5.22}\label{eq:5.5.22}
\end{equation}
By using either the eigenvalue conditions \eqref{eq:5.5.14} or the explicit formulas
\eqref{eq:5.5.17} and \eqref{eq:5.5.18}, we find at zero momentum:
\begin{equation}
N(\mathbf0)=\frac{\mathbf{1}+b_u\beta}{2(2\pi)^3},
\qquad
M(\mathbf0)=\frac{\mathbf{1}+b_v\beta}{2(2\pi)^3}.
\tag{5.5.23}\label{eq:5.5.23}
\end{equation}
From Eqs.~\eqref{eq:5.5.6} and \eqref{eq:5.5.7} we have then
\begin{equation}
N(\mathbf p)=\frac{m}{2p^0}D(L(\mathbf p))
[\mathbf{1}+b_u\beta]D^\dagger(L(\mathbf p)),
\tag{5.5.24}\label{eq:5.5.24}
\end{equation}
\begin{equation}
M(\mathbf p)=\frac{m}{2p^0}D(L(\mathbf p))
[\mathbf{1}+b_v\beta]D^\dagger(L(\mathbf p)).
\tag{5.5.25}\label{eq:5.5.25}
\end{equation}
The pseudounitarity condition \eqref{eq:5.4.32} yields
\[
D(L(\mathbf p))\beta D^\dagger(L(\mathbf p))=\beta
\]
and
\[
D(L(\mathbf p))D^\dagger(L(\mathbf p))
=D(L(\mathbf p))\beta D^{-1}(L(\mathbf p))\beta.
\]
We also recall that $\beta=-i\gamma^0$, so by using the Lorentz transformation
rule \eqref{eq:5.4.8} we have
\begin{equation}
D(L(\mathbf p))\beta D^{-1}(L(\mathbf p))
=-iL_\mu{}^0(\mathbf p)\gamma^\mu=-ip^\mu\gamma_\mu/m.
\tag{5.5.26}\label{eq:5.5.26}
\end{equation}
Putting this together, we find\footnote{Sometimes an extra factor
$\sqrt{p^0/m}$ is included in the Dirac spinors, so that $m$ appears in place
of $p^0$ in the denominators of the spin sums \eqref{eq:5.5.27} and \eqref{eq:5.5.28}. The
normalization convention used here has the advantage that it goes smoothly
over to the case $m=0$.}
\begin{equation}
N(\mathbf p)=\frac1{2p^0}[-ip^\mu\gamma_\mu+b_um]\beta,
\tag{5.5.27}\label{eq:5.5.27}
\end{equation}
\begin{equation}
M(\mathbf p)=\frac1{2p^0}[-ip^\mu\gamma_\mu+b_vm]\beta.
\tag{5.5.28}\label{eq:5.5.28}
\end{equation}
Using this in Eq.~\eqref{eq:5.5.20} yields finally
\begin{align}
[\psi_\ell(x),\psi_{\bar\ell}^\dagger(y)]_\mp
={}&\bigl(|\kappa|^2[-\gamma^\mu\partial_\mu+b_um]\beta\Delta_+(x-y)
\notag\\
&\mp|\lambda|^2[-\gamma^\mu\partial_\mu+b_vm]\beta\Delta_+(y-x)\bigr)_{\ell\bar\ell},
\tag{5.5.29}\label{eq:5.5.29}
\end{align}
where $\Delta_+$ is the function introduced in Section 5.2:
\[
\Delta_+(x)\equiv\int\frac{d^3p}{2p^0(2\pi)^3}e^{ip\cdot x}.
\]
We saw in Section 5.2 that for $x-y$ space-like $\Delta_+(x-y)$ is an even
function of $x-y$ and, of course, this implies that its first derivatives are
odd functions of $x-y$. Hence, in order that both the derivative and the
non-derivative terms in the commutator or anticommutator should vanish at
space-like separations, it is necessary and sufficient that
\begin{equation}
|\kappa|^2=\mp|\lambda|^2
\tag{5.5.30}\label{eq:5.5.30}
\end{equation}
and
\begin{equation}
|\kappa|^2b_u=\pm|\lambda|^2b_v.
\tag{5.5.31}\label{eq:5.5.31}
\end{equation}
Clearly Eq.~\eqref{eq:5.5.30} is only possible if we choose the bottom sign,
$\mp=+$; that is, \emph{the particles described by a Dirac field must be
fermions}. It is also then necessary that $|\kappa|^2=|\lambda|^2$ and
$b_u=-b_v$. Just as for scalars we have the freedom to redefine the relative
phase of the creation and annihilation operators to make the ratio
$\kappa/\lambda$ real, in which case $\kappa=\lambda$, and by adjusting the
overall scale and phase of the field $\psi$ we may then take
\begin{equation}
\kappa=\lambda=1.
\tag{5.5.32}\label{eq:5.5.32}
\end{equation}
Finally if we like we can replace $\psi$ with $\gamma^5\psi$, which changes the
sign of both $b_u$ and $b_v$, so we can always take
\begin{equation}
b_u=-b_v=+1.
\tag{5.5.33}\label{eq:5.5.33}
\end{equation}
For future use, we record here that the Dirac field is now
\begin{align}
\psi_\ell(x)=(2\pi)^{-3/2}\sum_\sigma\int d^3p\,
\bigl[u_\ell(\mathbf p,\sigma)e^{ip\cdot x}a_{\mathbf p,\sigma}
+v_\ell(\mathbf p,\sigma)e^{-ip\cdot x}a^{c\dagger}_{\mathbf p,\sigma}\bigr]
\tag{5.5.34}\label{eq:5.5.34}
\end{align}
while the coefficient functions at zero momentum are
\begin{equation}
u(\mathbf0,\tfrac12)=\frac1{\sqrt2}\begin{bmatrix}1\\0\\1\\0\end{bmatrix},
\qquad
u(\mathbf0,-\tfrac12)=\frac1{\sqrt2}\begin{bmatrix}0\\1\\0\\1\end{bmatrix},
\tag{5.5.35}\label{eq:5.5.35}
\end{equation}
\begin{equation}
v(\mathbf0,\tfrac12)=\frac1{\sqrt2}\begin{bmatrix}0\\1\\0\\-1\end{bmatrix},
\qquad
v(\mathbf0,-\tfrac12)=\frac1{\sqrt2}\begin{bmatrix}-1\\0\\1\\0\end{bmatrix}.
\tag{5.5.36}\label{eq:5.5.36}
\end{equation}
The spin sums are
\begin{equation}
N(\mathbf p)=\frac1{2p^0}[-ip^\mu\gamma_\mu+m]\beta,
\tag{5.5.37}\label{eq:5.5.37}
\end{equation}
\begin{equation}
M(\mathbf p)=\frac1{2p^0}[-ip^\mu\gamma_\mu-m]\beta,
\tag{5.5.38}\label{eq:5.5.38}
\end{equation}
so the anticommutator is given by Eq.~\eqref{eq:5.5.20} as
\begin{equation}
[\psi_\ell(x),\psi_{\bar\ell}^\dagger(y)]_+
=\{[-\gamma^\mu\partial_\mu+m]\beta\}_{\ell\bar\ell}\Delta(x-y).
\tag{5.5.39}\label{eq:5.5.39}
\end{equation}

Now let's return to the requirement that under a space inversion the field
$\psi(x)$ must transform into something proportional to $\psi(\mathcal P x)$.
For this to be possible the phases in Eqs.~\eqref{eq:5.5.15} and \eqref{eq:5.5.16} must be
equal, and so the intrinsic parities of particles and their antiparticles must
be related by
\begin{equation}
\eta^c=-\eta^*.
\tag{5.5.40}\label{eq:5.5.40}
\end{equation}
That is, \emph{the intrinsic parity $\eta\eta^c$ of a state consisting of a
spin $\frac12$ particle and its antiparticle is odd}. It is for this reason
that negative parity mesons like the $\rho^0$ and $J/\psi$ can be interpreted
as s-wave bound states of quark--antiquark pairs. Eqs.~\eqref{eq:5.5.15} and \eqref{eq:5.5.16}
now give the transformation of the causal Dirac field under space inversion as
\begin{equation}
P\psi(x)P^{-1}=\eta^*\beta\psi(\mathcal P x).
\tag{5.5.41}\label{eq:5.5.41}
\end{equation}

Before going on to the other inversions, this is a good place to mention that
Eqs.~\eqref{eq:5.5.14}, \eqref{eq:5.5.33}, and \eqref{eq:5.5.26} show that $u(\mathbf p,\sigma)$ and
$v(\mathbf p,\sigma)$ are eigenvectors of $-ip^\mu\gamma_\mu/m$ with
eigenvalues $+1$ and $-1$, respectively:
\begin{equation}
(ip^\mu\gamma_\mu+m)u(\mathbf p,\sigma)=0,
\qquad
(-ip^\mu\gamma_\mu+m)v(\mathbf p,\sigma)=0.
\tag{5.5.42}\label{eq:5.5.42}
\end{equation}
It follows then that the field \eqref{eq:5.5.33} satisfies the differential equation
\begin{equation}
(\gamma^\mu\partial_\mu+m)\psi(x)=0.
\tag{5.5.43}\label{eq:5.5.43}
\end{equation}
This is the celebrated Dirac equation for a free particle of spin $\frac12$.
From the point of view adopted here, the free-particle Dirac equation is
nothing but a Lorentz-invariant record of the convention that we have used in
putting together the two irreducible representations of the proper
orthochronous Lorentz group to form a field that transforms simply also under
space inversion.

In order to work out the charge-conjugation and time-reversal properties of
the Dirac field, we will need expressions for the complex-conjugates of the
$u$ and $v$ coefficient functions. These functions are real for zero momentum,
but to obtain the coefficient functions at finite momentum we have to multiply
with the complex matrix $D(L(\mathbf p))$. From Eq.~\eqref{eq:5.4.41} we see that for
general real $\omega_{\mu\nu}$:
\[
[\exp(\tfrac12i\mathcal{J}^{\mu\nu}\omega_{\mu\nu})]^*
=\beta\mathcal{C}\exp(\tfrac12i\mathcal{J}^{\mu\nu}\omega_{\mu\nu})
\mathcal{C}^{-1}\beta
\]
and so in particular
\[
D(L(\mathbf p))^*=\beta\mathcal{C}D(L(\mathbf p))\mathcal{C}^{-1}\beta.
\]
We also note that $\mathcal{C}^{-1}\beta u(\mathbf0,\sigma)=-v(\mathbf0,\sigma)$
and $\mathcal{C}^{-1}\beta v(\mathbf0,\sigma)=-u(\mathbf0,\sigma)$, so
\begin{equation}
u_\ell^*(\mathbf p,\sigma)=-(\beta\mathcal{C}v(\mathbf p,\sigma))_\ell,
\tag{5.5.44}\label{eq:5.5.44}
\end{equation}
\begin{equation}
v_\ell^*(\mathbf p,\sigma)=-(\beta\mathcal{C}u(\mathbf p,\sigma))_\ell.
\tag{5.5.45}\label{eq:5.5.45}
\end{equation}
In order for the field to transform under charge-conjugation into another
field with which it commutes at space-like separations, it is necessary again
that the charge conjugation parities of the particle and antiparticle be
related by
\begin{equation}
\xi^c=\xi^*.
\tag{5.5.46}\label{eq:5.5.46}
\end{equation}
In this case, the field transforms as
\begin{equation}
C\psi(x)C^{-1}=-\xi^*\beta\mathcal{C}\psi^*(x).
\tag{5.5.47}\label{eq:5.5.47}
\end{equation}
(We are calling the Hermitian adjoint of the field on the right-hand side
$\psi^*$ instead of $\psi^\dagger$ to emphasize that this is still a column
vector, not a row.)

Although we have been distinguishing particles from their antiparticles, we
have not ruled out the possibility that the two are actually identical. Such
spin $\frac12$ particles are called \emph{Majorana fermions}. Following the
same reasoning that led to Eq.~\eqref{eq:5.5.47}, the Dirac field of such a particle
must satisfy the reality condition
\begin{equation}
\psi(x)=-\beta\mathcal{C}\psi^*(x).
\tag{5.5.48}\label{eq:5.5.48}
\end{equation}
For Majorana fermions the intrinsic space-inversion parity must be imaginary,
$\eta=\mathbin{\pm}i$, while the charge-conjugation parity must be real,
$\xi=\mathbin{\pm}1$.

There is an important difference between fermions and bosons in the intrinsic
charge-conjugation phase of states consisting of a particle and its
antiparticle. Such a state may be written
\[
\ket{\Phi}=\sum_{\sigma,\sigma'}\int d^3p\int d^3p'\,
\chi(\mathbf p,\sigma;\mathbf p',\sigma')a^\dagger_{\mathbf p,\sigma}
a^{c\dagger}_{\mathbf p',\sigma'}\ket{\Omega},
\]
where $\ket{\Omega}$ is the vacuum state. Under charge-conjugation, this state
is transformed into
\[
C\ket{\Phi}=\xi^*\xi^{c*}\sum_{\sigma,\sigma'}\int d^3p\int d^3p'\,
\chi(\mathbf p,\sigma;\mathbf p',\sigma')a^{c\dagger}_{\mathbf p,\sigma}
a^\dagger_{\mathbf p',\sigma'}\ket{\Omega}.
\]
Interchanging the variables of integration and summation and using the
anticommutation of the creation operators and Eq.~\eqref{eq:5.5.46}, we can rewrite
this as
\[
C\ket{\Phi}=-\sum_{\sigma,\sigma'}\int d^3p\int d^3p'\,
\chi(\mathbf p',\sigma';\mathbf p,\sigma)a^\dagger_{\mathbf p,\sigma}
a^{c\dagger}_{\mathbf p',\sigma'}\ket{\Omega}.
\]
That is, \emph{the intrinsic charge-conjugation parity of a state consisting
of a particle described by a Dirac field and its antiparticle is odd}, in the
sense that if the wave function $\chi$ of the state is even or odd under
interchange of the momenta and spins of the particle and antiparticle, then the
charge-conjugation operator applied to such a state gives a sign $-1$ or $+1$,
respectively. The classic example here is positronium, the bound state of an
electron and a positron. The two lowest states are a pair of nearly degenerate
s-wave states with total spin $s=0$ and $s=1$, known respectively as para- and
ortho-positronium. The wave function for these two states is even under
interchange of momenta and odd or even respectively under interchange of spin
$z$-components, so para- and ortho-positronium have $C=+1$ and $C=-1$,
respectively. These values are dramatically confirmed in the decay modes of
positronium: para-positronium decays rapidly into a pair of photons (each of
which has $C=-1$), while ortho-positronium can only decay much more slowly into
three or more photons. In the same way, single $\rho^0$ and $\omega^0$ mesons
are produced as resonances in high-energy electron--positron annihilation
through a one-photon intermediate state, so they must have $C=-1$, which is
consistent with their interpretation as quark--antiquark bound states with
orbital angular momentum zero and total quark spin one.

Now we come to time-reversal. Recall the transformation properties of the
particle annihilation and antiparticle creation operators given by Eq.
\eqref{eq:4.2.15}:
\begin{equation}
T a_{\mathbf p,\sigma}T^{-1}
=\zeta^*(-1)^{1/2-\sigma}a_{-\mathbf p,-\sigma},
\tag{5.5.49}\label{eq:5.5.49}
\end{equation}
\begin{equation}
T a^{c\dagger}_{\mathbf p,\sigma}T^{-1}
=\zeta^c(-1)^{1/2-\sigma}a^{c\dagger}_{-\mathbf p,-\sigma}.
\tag{5.5.50}\label{eq:5.5.50}
\end{equation}
Time-reversal of the field \eqref{eq:5.5.34} thus gives
\begin{align*}
T\psi_\ell(x)T^{-1}={}&(2\pi)^{-3/2}\sum_\sigma\int d^3p\,
(-1)^{1/2-\sigma}\\
&\cdot[\zeta^*u_\ell^*(\mathbf p,\sigma)e^{-ip\cdot x}a_{-\mathbf p,-\sigma}
+\zeta^c v_\ell^*(\mathbf p,\sigma)e^{ip\cdot x}
a^{c\dagger}_{-\mathbf p,-\sigma}].
\end{align*}
In order to put this back in the form given for $\psi$, we shall redefine the
variables of integration and summation as $-\mathbf p$ and $-\sigma$, so we
need formulas for $u^*(-\mathbf p,-\sigma)$ and
$v^*(-\mathbf p,-\sigma)$ in terms of $u(\mathbf p,\sigma)$ and
$v(\mathbf p,\sigma)$, respectively. For this purpose, we can use the fact
that $\mathcal{J}^{i0}$ anticommutes with $\beta$ and commutes with $\gamma^5$
together with our former result for $D(L(\mathbf p))^*$ to write
\[
D^*(L(-\mathbf p))=\gamma^5\beta D^*(L(\mathbf p))\beta\gamma^5
=\gamma^5\mathcal{C}D(L(\mathbf p))\mathcal{C}^{-1}\gamma^5.
\]
Also, Eqs.~\eqref{eq:5.4.36} and \eqref{eq:5.5.35}--\eqref{eq:5.5.36} give
\[
\gamma^5\mathcal{C}^{-1}u(\mathbf0,-\sigma)
=(-1)^{1/2-\sigma}u(\mathbf0,\sigma),
\]
\[
\gamma^5\mathcal{C}^{-1}v(\mathbf0,-\sigma)
=(-1)^{1/2-\sigma}v(\mathbf0,\sigma),
\]
so
\begin{equation}
(-1)^{1/2+\sigma}u^*(-\mathbf p,-\sigma)
=-\gamma^5\mathcal{C}u(\mathbf p,\sigma),
\tag{5.5.51}\label{eq:5.5.51}
\end{equation}
\begin{equation}
(-1)^{1/2+\sigma}v^*(-\mathbf p,-\sigma)
=-\gamma^5\mathcal{C}v(\mathbf p,\sigma).
\tag{5.5.52}\label{eq:5.5.52}
\end{equation}
We see then that in order for time-reversal to take the Dirac field into
something proportional to itself at the time-reversed point (with which it
would anticommute at space-like separations) it is necessary that the intrinsic
time-reversal phases be related by
\begin{equation}
\zeta^c=\zeta^*
\tag{5.5.53}\label{eq:5.5.53}
\end{equation}
and in this case
\begin{equation}
T\psi(x)T^{-1}=-\zeta^*\gamma^5\mathcal{C}\psi(-\mathcal T x).
\tag{5.5.54}\label{eq:5.5.54}
\end{equation}

Now let us consider how to construct scalar interaction densities out of the
Dirac fields and their adjoints. As already mentioned the Dirac representation
is not unitary, so $\psi^\dagger\psi$ is not a scalar. To deal with this
complication it is convenient to define a new sort of adjoint:
\begin{equation}
\bar\psi\equiv\psi^\dagger\beta.
\tag{5.5.55}\label{eq:5.5.55}
\end{equation}
Using the pseudounitarity condition \eqref{eq:5.4.32}, we see that the fermion bilinears
constructed with $\bar\psi$ have the Lorentz transformation property
\begin{equation}
U_0(\Lambda)[\bar\psi(x)M\psi(x)]U_0^{-1}(\Lambda)
=\bar\psi(\Lambda x)D^{-1}(\Lambda)MD(\Lambda)\psi(\Lambda x).
\tag{5.5.56}\label{eq:5.5.56}
\end{equation}
Also, under a space inversion
\begin{equation}
P[\bar\psi(x)M\psi(x)]P^{-1}
=\bar\psi(\mathcal P x)\beta M\beta\psi(\mathcal P x).
\tag{5.5.57}\label{eq:5.5.57}
\end{equation}
Taking the matrix $M$ as $\mathbf{1}$, $\gamma^\mu$,
$\mathcal{J}^{\mu\nu}$, $\gamma^5\gamma^\mu$, or $\gamma^5$ yields a bilinear
$\bar\psi M\psi$ that transforms as a scalar, vector, tensor, axial vector,
and pseudoscalar, respectively. (The terms `axial' and `pseudo' indicate that
these have space-inversion properties opposite to those of ordinary vectors
and scalars: a pseudoscalar has negative parity, while the space and time
components of an axial vector have positive and negative parity, respectively.)
These results apply also when the two fermion fields in the bilinear refer to
different particle species, except that in this case a space inversion also
yields a ratio of the intrinsic parities.

For instance, the original Fermi theory of beta decay involved an interaction
density proportional to $\bar\psi_p\gamma^\mu\psi_n\,
\bar\psi_e\gamma_\mu\psi_\nu$. Later it was realized that the most general
Lorentz-invariant and parity-conserving non-derivative beta decay interaction
takes the form of a linear combination of products like this, with $\gamma_\mu$
replaced with any one of the five covariant types of $4\cdot4$ matrices
$\mathbf{1}$, $\gamma^\mu$, $\mathcal{J}^{\mu\nu}$,
$\gamma^5\gamma^\mu$, or $\gamma^5$. (As discussed in Chapter 2, we are
defining the space-inversion operator so that the proton, neutron, and electron
all have parity $+1$. If the neutrino is massless then its parity may also be
defined as $+1$, if necessary by replacing the neutrino field with
$\gamma^5\psi_\nu$.) When Lee and Yang~\hyperref[ref:5.7]{[7]} called parity conservation into
question in 1956, they expanded the list of possible non-derivative
interactions to include ten terms proportional to
$\bar\psi_pM\psi_n\,\bar\psi_eM\psi_\nu$ and also
$\bar\psi_pM\psi_n\,\bar\psi_eM\gamma^5\psi_\nu$, with $M$ running over the
matrices $\mathbf{1}$, $\gamma^\mu$, $\mathcal{J}^{\mu\nu}$,
$\gamma^5\gamma^\mu$, or $\gamma^5$.

It is also of some interest to study the charge-conjugation properties of these
bilinears. Using Eqs.~\eqref{eq:5.5.47} and \eqref{eq:5.4.35}--\eqref{eq:5.4.39}, we have
\begin{align}
C(\bar\psi M\psi)C^{-1}
&=(\beta\mathcal{C}\psi)^{\mathrm T}\beta M(\beta\mathcal{C}\psi^*)
=-(\beta\mathcal{C}\psi^*)^{\mathrm T}M^{\mathrm T}\mathcal{C}\psi
\notag\\
&=\bar\psi\mathcal{C}^{-1}M^{\mathrm T}\mathcal{C}\psi
=\mathbin{\pm}\bar\psi M\psi,
\tag{5.5.58}\label{eq:5.5.58}
\end{align}
the sign in the last expression being $+$ for the matrices $\mathbf{1}$,
$\gamma^5\gamma_\mu$, and $\gamma^5$, and $-$ for $\gamma_\mu$ and
$\mathcal{J}_{\mu\nu}$. (The minus sign in the first line arises from Fermi
statistics. We ignore a c-number anticommutator.) A boson field that interacts
with the current $\bar\psi M\psi$ must therefore have $C=+1$ for scalars,
pseudoscalars, or axial vectors, and $C=-1$ for vectors or antisymmetric
tensors. This is one way of seeing that the $\pi^0$ (which couples to
pseudoscalar or axial-vector nucleon currents) has $C=+1$, while the photon
has $C=-1$.
